\documentclass[a4paper, serif, 10pt]{moderncv}

%% moderncv
% color and style
\moderncvcolor{black}
\moderncvstyle{banking}

%% page layout/margins
\usepackage[top=2.5cm, bottom=2.5cm, left=1.5cm, right=1.5cm]{geometry}

\nopagenumbers{}

%% Babel config for Indonesian (Bahasa Indonesia) language
% ref:
% - https://latex3.github.io/babel/guides/locale-indonesian.html
\usepackage[indonesian]{babel}

%% fontspec
\usepackage{fontspec}

\IfFontExistsTF{Linux Libertine}{
  \setmainfont[Ligatures={TeX, Common}, Numbers=OldStyle]{Linux Libertine}
}{}
\IfFontExistsTF{Linux Libertine O}{
  \setmainfont[Ligatures={TeX, Common}, Numbers=OldStyle]{Linux Libertine O}
}{}

%% CTeX
% CJK support, used to show CJK characters in the resume
%
% - heading=false: disable ctex heading
% - punct=kaiming: use kaiming punctuations styles for CJK
% - scheme=plain: use plain scheme, do not override \normalsize font size
% - space=auto: space settings for CJK characters
%
% CJK fonts are resolved per locale language, e.g. Noto Serif/Sans CJK JP
% for Japanese resumes, with fallback to the next region if not installed
%
% ref:
% - http://ctan.mirrorcatalogs.com/language/chinese/ctex/ctex.pdf
\usepackage[UTF8, heading=false, punct=kaiming, scheme=plain, space=auto]{ctex}

\IfFontExistsTF{Noto Serif CJK SC}{
  \setCJKmainfont{Noto Serif CJK SC}
}{}
\IfFontExistsTF{Noto Sans CJK SC}{
  \setCJKsansfont{Noto Sans CJK SC}
}{}

%% line spacing
\usepackage{setspace}
\setstretch{1.125}

%% URL styling - use normal text instead of monospace
\urlstyle{same}

%% Patch \httplink and \httpslink to handle full URLs
% The original moderncv macros always prepend http:// or https:// to URLs.
% This causes issues when the URL already has a protocol (e.g., https://example.com)
% resulting in malformed URLs like https://https://example.com.
% This patch checks if the URL already contains "://" and uses it directly if so.
% Note: We use \str_set:Nx (x-expansion) to expand macros like \@homepage first.
\ExplSyntaxOn
\str_new:N \g_yamlresume_url_str

\RenewDocumentCommand{\httpslink}{O{}m}{
  \str_set:Nx \g_yamlresume_url_str {#2}
  \str_if_in:NnTF \g_yamlresume_url_str {://}
    { \url{#2} }
    { \url{https://#2} }
}

\RenewDocumentCommand{\httplink}{O{}m}{
  \str_set:Nx \g_yamlresume_url_str {#2}
  \str_if_in:NnTF \g_yamlresume_url_str {://}
    { \url{#2} }
    { \url{http://#2} }
}
\ExplSyntaxOff

\name{Aditya Pratama Nugroho}{}
\title{Insinyur AI | Pembelajaran Mesin \& MLOps}
\phone[mobile]{+62 812-3456-7890}
\email{aditya.pratama.ai@gmail.com}
\homepage{https://adityapratama.id}

\address{Jl. Gatot Subroto No. 88, Jakarta Selatan, DKI Jakarta, Indonesia, 12930}{}{}

\extrainfo{{\small \faLinkedin}\ \href{https://www.linkedin.com/in/adityapratamanugroho}{@adityapratamanugroho} {} {} {} • {} {} {} 
{\small \faGithub}\ \href{https://github.com/adityapratama}{@adityapratama} {} {} {} • {} {} {} 
{\small \faStackOverflow}\ \href{https://stackoverflow.com/users/1234567/aditya-pratama}{@aditya-pratama}}

\begin{document}

\maketitle

\section{Informasi Dasar}

\cvline{}{Insinyur AI dengan pengalaman 6+ tahun merancang, membangun, dan menerapkan solusi machine learning serta deep learning untuk skala produksi. Terbiasa menangani siklus penuh mulai dari eksplorasi data, feature engineering, pelatihan model, evaluasi, hingga deployment dan monitoring di cloud.
%
\begin{itemize}
\item Menguasai Python, PyTorch, TensorFlow, dan scikit-learn untuk pengembangan model prediktif, NLP, dan computer vision.
\item Berpengalaman membangun pipeline MLOps menggunakan Docker, Kubernetes, MLflow, dan Airflow.
\item Kolaboratif, berorientasi hasil, dan mampu menerjemahkan kebutuhan bisnis menjadi solusi AI yang terukur.
\end{itemize}}
\section{Pendidikan}

\cventry{Agu 2014–Jul 2018}
        {Sarjana, Ilmu Komputer, Nilai: 3.72}
        {Universitas Indonesia}
        {\url{https://www.ui.ac.id}}
        {}
        {\begin{itemize}
\item Lulus dengan predikat cum laude dan fokus pada bidang kecerdasan buatan serta pembelajaran mesin.
\item Menjadi asisten dosen untuk mata kuliah Pembelajaran Mesin dan Struktur Data.
\item Menyelesaikan skripsi tentang klasifikasi teks berita berbahasa Indonesia menggunakan model transformer.
\end{itemize}
\textbf{Mata Kuliah}: Struktur Data dan Algoritma, Pembelajaran Mesin, Pengolahan Bahasa Alami, Visi Komputer, Basis Data, Sistem Terdistribusi, Statistika Komputasi, Kecerdasan Buatan}
\section{Pengalaman Kerja}

\cventry{Jan 2022–Sekarang}
        {Insinyur AI Senior}
        {PT Nusantara AI Solutions}
        {\url{https://www.nusantara-ai.co.id}}
        {}
        {\begin{itemize}
\item Memimpin tim 5 engineer dalam pengembangan platform rekomendasi berbasis deep learning yang melayani 2 juta pengguna aktif.
\item Merancang arsitektur model NLP untuk analisis sentimen dan ekstraksi entitas dengan akurasi F1 di atas 0,92.
\item Membangun pipeline MLOps end-to-end menggunakan MLflow, Airflow, Docker, dan Kubernetes untuk 30+ model produksi.
\item Mengoptimalkan biaya inferensi sebesar 35\% melalui model quantization, caching, dan autoscaling.
\item Berkolaborasi dengan product manager dan data analyst untuk menerjemahkan kebutuhan bisnis menjadi roadmap AI.
\end{itemize}
\textbf{Kata Kunci}: Deep Learning, NLP, MLOps, Kubernetes, Leadership}

\cventry{Mar 2019–Des 2021}
        {Insinyur Pembelajaran Mesin}
        {PT Data Cerdas Indonesia}
        {\url{https://www.datacerdas.co.id}}
        {}
        {\begin{itemize}
\item Mengembangkan model prediksi churn dan credit scoring untuk industri finansial dengan AUC 0,88.
\item Membangun fitur dan pipeline data menggunakan Python, Spark, dan SQL untuk dataset 50 juta transaksi.
\item Menerapkan model ke layanan API menggunakan FastAPI dan Docker dengan latensi rata-rata di bawah 80 ms.
\item Melakukan A/B testing dan monitoring model untuk memastikan performa dan stabilitas di lingkungan produksi.
\end{itemize}
\textbf{Kata Kunci}: Machine Learning, Spark, Python, FastAPI, Credit Scoring}

\cventry{Agu 2018–Feb 2019}
        {Ilmuwan Data}
        {PT Solusi Analitik Global}
        {\url{https://www.solusi-analitik.co.id}}
        {}
        {\begin{itemize}
\item Menganalisis data pelanggan untuk mengidentifikasi peluang cross-selling dan segmentasi pasar.
\item Membuat dashboard interaktif dan laporan insight menggunakan Tableau dan Python.
\item Berkolaborasi dengan tim penjualan untuk merancang kampanye berbasis data yang meningkatkan konversi 12\%.
\end{itemize}
\textbf{Kata Kunci}: Data Analysis, Tableau, Segmentation, Python}
\section{Bahasa}

\cvline{Indonesia}{Bahasa Ibu / Bilingual \hfill \textbf{Kata Kunci}: Bahasa ibu}
\cvline{Inggris}{Tingkat Profesional Penuh \hfill \textbf{Kata Kunci}: TOEFL iBT 105, IELTS 7.5}
\section{Keahlian}

\cvline{Pembelajaran Mesin}{Ahli \hfill \textbf{Kata Kunci}: Python, PyTorch, TensorFlow, scikit-learn, XGBoost}
\cvline{Pembelajaran Mendalam}{Mahir \hfill \textbf{Kata Kunci}: NLP, Computer Vision, Transformers, CNN, RNN}
\cvline{MLOps}{Mahir \hfill \textbf{Kata Kunci}: MLflow, Airflow, Docker, Kubernetes, CI/CD}
\cvline{Rekayasa Data}{Mahir \hfill \textbf{Kata Kunci}: SQL, Spark, Kafka, PostgreSQL, BigQuery}
\cvline{Komputasi Awan}{Menengah \hfill \textbf{Kata Kunci}: AWS, GCP, SageMaker, Vertex AI, Terraform}
\section{Penghargaan}

\cventry{Nov 2023}
        {PT Nusantara AI Solutions}
        {Penghargaan Inovasi AI Terbaik}
        {}
        {}
        {Diberikan atas keberhasilan memimpin pengembangan platform rekomendasi berbasis AI yang meningkatkan engagement pengguna sebesar 28\%.}

\cventry{Jun 2017}
        {Universitas Indonesia}
        {Dean's List}
        {}
        {}
        {Diberikan kepada mahasiswa dengan prestasi akademik tinggi dan konsisten selama dua semester berturut-turut.}
\section{Sertifikat}

\cventry{Mar 2023}
        {Amazon Web Services}
        {AWS Certified Machine Learning - Specialty}
        {\url{https://aws.amazon.com/certification/certified-machine-learning-specialty/}}
        {}
        {}

\cventry{Sep 2021}
        {Google}
        {TensorFlow Developer Certificate}
        {\url{https://www.tensorflow.org/certificate}}
        {}
        {}

\cventry{Feb 2022}
        {Google Cloud}
        {Professional Data Engineer}
        {\url{https://cloud.google.com/certification/data-engineer}}
        {}
        {}
\section{Publikasi}

\cventry{Okt 2021}
        {Klasifikasi Teks Berita Berbahasa Indonesia Menggunakan IndoBERT}
        {Jurnal Ilmu Komputer dan Informasi}
        {\url{https://jiki.cs.ui.ac.id/index.php/jiki}}
        {}
        {\begin{itemize}
\item Menerapkan model IndoBERT untuk klasifikasi topik berita berbahasa Indonesia.
\item Mencapai akurasi 94,2\% pada dataset 20.000 berita dari berbagai portal daring.
\item Menganalisis pengaruh fine-tuning dan ukuran data terhadap performa model.
\end{itemize}}

\cventry{Agu 2020}
        {Deteksi Anomali Transaksi Keuangan dengan Autoencoder}
        {Seminar Nasional Artificial Intelligence}
        {\url{https://semnas-ai.id/prosiding}}
        {}
        {\begin{itemize}
\item Mengembangkan model autoencoder untuk mendeteksi anomali pada data transaksi keuangan.
\item Membandingkan performa dengan Isolation Forest dan One-Class SVM.
\item Menunjukkan peningkatan recall sebesar 15\% pada data yang tidak seimbang.
\end{itemize}}
\section{Proyek}

\cventry{Jan 2023–Jun 2023}
        {Chatbot berbasis LLM untuk layanan pelanggan berbahasa Indonesia.}
        {Asisten Virtual Berbahasa Indonesia}
        {\url{https://github.com/adityapratama/asisten-virtual-id}}
        {}
        {\begin{itemize}
\item Membangun chatbot dengan arsitektur RAG menggunakan LangChain, vector database, dan model open-source.
\item Mengintegrasikan API ke kanal WhatsApp dan web dengan latensi respons di bawah 2 detik.
\item Menerapkan evaluasi otomatis untuk kualitas jawaban dan deteksi halusinasi.
\end{itemize}
\textbf{Kata Kunci}: LLM, RAG, LangChain, Chatbot}

\cventry{Mar 2020–Nov 2020}
        {Deteksi kerusakan jalan berbasis computer vision dari citra kendaraan.}
        {Sistem Deteksi Kerusakan Jalan}
        {\url{https://github.com/adityapratama/deteksi-jalan}}
        {}
        {\begin{itemize}
\item Melatih model YOLOv5 untuk mendeteksi retak dan lubang jalan pada citra dashcam.
\item Mencapai mAP 0,81 pada dataset 5.000 gambar yang dianotasi secara manual.
\item Mengemas model sebagai layanan inferensi menggunakan FastAPI dan Docker.
\end{itemize}
\textbf{Kata Kunci}: Computer Vision, YOLOv5, FastAPI}
\section{Minat}

\cvline{Teknologi}{Open Source, LLM, Edge AI}
\cvline{Olahraga}{Lari, Bersepeda, Bulu Tangkis}
\cvline{Musik}{Gitar, Jazz}
\section{Kegiatan Sukarela}

\cventry{Feb 2021–Sekarang}
        {Mentor Machine Learning}
        {Komunitas AI Indonesia}
        {\url{https://www.aiindonesia.org}}
        {}
        {\begin{itemize}
\item Menjadi mentor untuk program bootcamp machine learning bagi mahasiswa dan profesional muda.
\item Membantu peserta mengembangkan proyek akhir dan mempersiapkan portofolio AI.
\item Menyelenggarakan workshop tentang MLOps dan deployment model.
\end{itemize}}
    

\cventry{Sep 2016–Jun 2018}
        {Asisten Dosen}
        {Universitas Indonesia}
        {\url{https://www.ui.ac.id}}
        {}
        {\begin{itemize}
\item Membantu dosen dalam menyusun materi dan soal untuk mata kuliah Pembelajaran Mesin.
\item Membimbing mahasiswa dalam tugas pemrograman dan proyek kelompok.
\item Memberikan sesi konsultasi mingguan tentang konsep algoritma dan statistika.
\end{itemize}}
    
\end{document}